% Same records used by SQL example 10.2; the junction sits between its parents.
\begin{flushleft}
\begin{tikzpicture}[every node/.style={font=\small,inner sep=0pt,align=left}]
\node (student) {\begin{tabular}{|l|l|}\hline
\multicolumn{2}{|l|}{\textbf{Student}}\\\hline
\chSixHeader StudentID & Name\\\hline
101 & Rupa\\\hline 102 & Nabil\\\hline
\end{tabular}};
\node[anchor=north west] (enrolment) at ([xshift=12mm]student.north east)
{\begin{tabular}{|l|l|}\hline
\multicolumn{2}{|l|}{\textbf{Enrolment}}\\\hline
\chSixHeader StudentID & CourseID\\\hline
101 & C01\\\hline 101 & C02\\\hline 102 & C01\\\hline
\end{tabular}};
\node[anchor=north west] (course) at ([xshift=12mm]enrolment.north east)
{\begin{tabular}{|l|l|}\hline
\multicolumn{2}{|l|}{\textbf{Course}}\\\hline
\chSixHeader CourseID & CourseName\\\hline
C01 & ICT\\\hline C02 & English\\\hline
\end{tabular}};
\draw[bookarrow] (student.east)--node[pos=.2,above=2pt]{1}node[pos=.8,above=2pt]{M}(enrolment.west);
\draw[bookarrow] (course.west)--node[pos=.2,above=2pt]{1}node[pos=.8,above=2pt]{M}(enrolment.east);
\node[anchor=north west,font=\footnotesize] at ([yshift=-2mm]student.south west) {StudentID = PK};
\node[anchor=north west,font=\footnotesize,text width=38mm] at ([yshift=-2mm]enrolment.south west) {StudentID ও CourseID = FK};
\node[anchor=north west,font=\footnotesize] at ([yshift=-2mm]course.south west) {CourseID = PK};
\end{tikzpicture}
\captionof{table}{Student ও Course-এর মাঝখানে Enrolment Junction Table}
\end{flushleft}
